% !PRINTED 528
\weberpage{528}
Aus $P_2$ gehen ebenfalls drei conjugirte Gruppen hervor, die aber ausser 1 keinen gemeinschaftlichen Theiler haben.

\gesperrt{Es giebt also ausser $P$, $Q$, $P_1$, $Q_1$, $P_2$ und den mit $P_1$ und $P_2$ conjugirten Theilern von $P$ keine transitiven Permutationsgruppen von vier Elementen.}

Unter den intransitiven Gruppen verdient noch die Gruppe (3):
\[
P_3 = 1, \ (0,2), \ (1,3), \ (0,2)(1,3),
\]
hervorgehoben zu werden, die gleichfalls zu einem System von drei conjugirten Gruppen Anlass giebt.

\begin{center}
\S.\,161.

\gesperrt{Auflösung der biquadratischen Gleichungen.}
\end{center}

Die verschiedenen Methoden der Auflösung der biquadratischen Gleichungen unterscheiden sich von einander durch die Reihenfolge, in der die verschiedenen Divisoren der Gruppe $P$ benutzt werden.

Es seien $\alpha$, $\alpha_1$, $\alpha_2$, $\alpha_3$ die Wurzeln der biquadratischen Gleichung
\weq{1}{f(x) = x^4 - a_1x^3 + a_2x^2 - a_3x + a_4 = 0.}
Wir bezeichnen mit $D$ die Discriminante und setzen
\weq{2}{\sqrt{D} = (\alpha-\alpha_1)(\alpha-\alpha_2)(\alpha-\alpha_3)(\alpha_1-\alpha_2)(\alpha_1-\alpha_3)(\alpha_2-\alpha_3).}
Ausserdem erinnern wir noch an die Ausdrücke für die Invarianten (\S.\,64):
\weq{3}{A = a_2^2 - 3a_1a_3 + 12a_4}
\weq{4}{B = 27a_1^2a_4 + 27a_3^2 + 2a_2^3 - 72a_2a_4 - 9a_1a_2a_3,}
durch die man $D$ in der Form ausdrückt:
\weq{5}{27D = 4A^3 - B^2.}
Adjungirt man zunächst $\sqrt{D}$, so reducirt sich die Gruppe der Gleichung auf $Q$, und wenn man weiter eine zu $Q_1$ gehörige Function sucht, so wird durch diese die Gruppe auf $Q_1$ reducirt. Da $Q_1$ innerhalb $Q$ den Index 3 hat, so wird eine solche Function von einer cubischen Gleichung abhängen, und zwar, da $Q_1$ ein Normaltheiler von $Q$ ist, von einer Normalgleichung. Nehmen wir z.\,B.
\weq{6}{y = (\alpha-\alpha_1)(\alpha_2-\alpha_3)}
%CONT>>
